\documentclass[11pt,a4paper]{article}

\usepackage{ctex}
\usepackage[T1]{fontenc}
\usepackage[top=2.5cm, bottom=2.5cm, left=2.5cm, right=2.5cm]{geometry}
\usepackage{booktabs}
\usepackage{array}
\usepackage{enumitem}
\usepackage{tabularx}
\usepackage{amssymb}
\usepackage{xcolor}
\definecolor{primary}{HTML}{1a1a2e}
\definecolor{accent}{HTML}{3b82f6}
\definecolor{success}{HTML}{22c55e}
\definecolor{danger}{HTML}{ef4444}
\definecolor{warning}{HTML}{f59e0b}
\definecolor{graytext}{HTML}{6b7280}
\usepackage[hidelinks]{hyperref}
\hypersetup{colorlinks=true, linkcolor=primary, urlcolor=accent}
\usepackage{titlesec}
\titleformat{\section}{\Large\bfseries\color{primary}}{}{0em}{}[\titlerule]
\titleformat{\subsection}{\large\bfseries\color{primary}}{}{0em}{}
\usepackage{fancyhdr}
\pagestyle{fancy}
\fancyhf{}
\fancyhead[L]{\small\color{graytext} 企业级 AI Agent 应用商店}
\fancyhead[R]{\small\color{graytext} RFC}
\fancyfoot[C]{\small\color{graytext} \thepage}
\renewcommand{\headrulewidth}{0.4pt}

\title{\textbf{RFC：请求意见/技术提案}\\[0.5em]\large Request for Comments}
\author{万能产品小助手（PM AI）\\魏子政 审阅}
\date{2026-05-06 / 版本 v1.0}

\begin{document}

\maketitle

\begin{center}
\begin{tabular}{>{\bfseries}p{3cm} p{10cm}}
\toprule
文档类型 & RFC（Request for Comments） \\
定位 & 请求意见/技术提案 -- 技术方案发出来让大家看看，有没有坑 \\
核心问题 & 行不行得通？有没有更好的选择？ \\
前置文档 & MRD（市场需求文档） \\
后续文档 & TRD（技术需求文档） \\
\bottomrule
\end{tabular}
\end{center}

\vspace{1em}

\section{技术栈（已锁定）}

\begin{table}[h]
\centering
\begin{tabularx}{\linewidth}{lXl}
\toprule
\textbf{层} & \textbf{技术选型} & \textbf{选择理由} \\
\midrule
前端框架 & Vue 3 + Element Plus & 国内部署友好，组件丰富，黑白灰极简风格实现成本低 \\
构建工具 & Vite & HMR 快，开发体验好，生态成熟 \\
状态管理 & Pinia & Vue 3 官方推荐，TypeScript 友好 \\
后端框架 & FastAPI & Python 生态，async 原生，Swagger 自动生成 \\
ORM & SQLAlchemy 2.0 (async) & 功能全面，async 支持，Alembic 迁移集成 \\
数据库 & PostgreSQL 16 & JSONB 灵活，关系查询强，企业级成熟度 \\
搜索引擎 & OpenSearch 2.x & 全文+语义搜索，中文分词（IK） \\
对象存储 & MinIO & S3 兼容，私有部署，轻量 \\
认证 & JWT (python-jose + passlib) & 无状态，适合 API 服务，无需额外组件 \\
邮箱验证 & 内存验证码（MVP） & MVP 阶段简化方案，预留 SMTP 接入点 \\
病毒扫描 & ClamAV（Docker） & 开源，clamav/clamav:latest 官方镜像 \\
静态分析 & Semgrep & 开源，规则丰富，CI/CD 友好 \\
LLM 分析 & qwen-turbo（DashScope） & 国内部署友好，API 调用简单 \\
容器编排 & Docker Compose & 开发/测试/部署统一，轻量够用 \\
\bottomrule
\end{tabularx}
\end{table}

\section{架构方案（已实现）}

\subsection{前端架构}

SPA 单页应用。Vue Router 路由守卫做权限控制（requiresAuth / requiresAdmin）。API 调用统一封装 axios 实例（api/client.ts），Pinia 管理认证状态（stores/auth.ts）。

\subsection{后端架构}

六层分层架构：
\begin{itemize}[nosep]
  \item \texttt{routers/} — API 路由层（auth / skills / reviews）
  \item \texttt{services/} — 业务逻辑层（scan\_service / skill\_package\_service）
  \item \texttt{models/} — ORM 模型层（User / Skill / ScanResult / Review）
  \item \texttt{schemas/} — Pydantic 数据模型层（请求/响应格式定义）
  \item \texttt{tasks/} — 后台任务层（扫描任务）
  \item \texttt{utils/} — 工具层（minio\_client / password）
\end{itemize}

全部 async/await，BackgroundTasks 驱动扫描。

\subsection{安全扫描管道}

上传 zip $\to$ Semgrep（静态代码分析）$\to$ ClamAV（病毒扫描）$\to$ qwen-turbo（LLM 语义分析）$\to$ 三层结果汇总 $\to$ 驱动审批。

\section{讨论点与决策}

\begin{table}[h]
\centering
\begin{tabularx}{\linewidth}{lp{2cm}X}
\toprule
\textbf{编号} & \textbf{状态} & \textbf{讨论内容与决策} \\
\midrule
R-001 & \textcolor{success}{已决定} & 前端选 Vue 3 + Element Plus（部署友好 + 黑白灰风格实现简单） \\
R-002 & \textcolor{success}{已决定} & 后端选 FastAPI（轻量 + Python 生态 + async 原生） \\
R-003 & \textcolor{success}{已决定} & 认证选 JWT（无状态，减少组件依赖） \\
R-004 & \textcolor{success}{已决定} & 编排选 Docker Compose（开发阶段够用） \\
R-005 & \textcolor{success}{已决定} & LLM 用 qwen-turbo + DashScope API（国内友好，成本低） \\
R-006 & \textcolor{success}{已决定} & 邮箱验证 MVP 用内存验证码，预留 SMTP 接入 \\
R-007 & \textcolor{success}{已决定} & 暂不需要消息队列，asyncio BackgroundTasks 够用 \\
R-008 & \textcolor{success}{已决定} & Skill 安装到 $\sim$/.openclaw/skills/，支持中文目录名 \\
R-009 & \textcolor{success}{已决定} & Skill 状态流转：scanning/pending\_review/published/offline/rejected \\
R-010 & \textcolor{success}{已决定} & 全部用户默认 admin 角色（MVP 简化） \\
\bottomrule
\end{tabularx}
\end{table}

\section{风险点}

\begin{itemize}[nosep]
  \item \textbf{LLM 扫描延迟}：qwen-turbo API 调用可能成为审批流程瓶颈（当前可接受）
  \item \textbf{OpenSearch 资源}：单节点内存占用 512MB-1GB（当前未启用全文搜索）
  \item \textbf{邮箱验证 MVP}：内存存储验证码，服务重启后丢失（后续接 SMTP）
  \item \textbf{ClamAV Docker}：首次扫描需更新病毒库，可能较慢
\end{itemize}

\section*{更新记录}

\begin{tabular}{lll}
\toprule
版本 & 日期 & 变更内容 \\
\midrule
v1.0 & 2026-05-06 & 全面重写：R-005/R-006 已决定（qwen-turbo + 内存验证码），新增 R-008~R-010（安装路径/状态流转/全管理员），补充实际架构描述 \\
v0.2 & 2026-05-06 & R-005/R-006 改为待拍板 \\
v0.1 & 2026-05-06 & 初始版本 \\
\bottomrule
\end{tabular}

\end{document}
